% !TeX program = pdflatex
% !BIB program = biber

\documentclass[11pt]{article}

\usepackage[a4paper,margin=1in]{geometry}
\usepackage[T1]{fontenc}
\usepackage[utf8]{inputenc}
\usepackage{lmodern}
\usepackage{amsmath,amssymb,amsthm,bm,mathtools}
\usepackage{enumitem}
\usepackage{array}
\usepackage{longtable}
\usepackage{booktabs}
\usepackage{xurl}
\usepackage[expansion=false]{microtype}
\setlength{\emergencystretch}{3em}
\usepackage{xcolor}
\usepackage{hyperref}
\usepackage[nameinlink,noabbrev]{cleveref}

\usepackage[backend=biber,style=numeric,sorting=none,maxnames=6,uniquename=false,url=false]{biblatex}
\addbibresource{cmotion_foundational_general.bib}
\addbibresource{cmotion_primary_research.bib}
\addbibresource{cmotion_internal_manuscripts.bib}

\AtEveryBibitem{\clearfield{note}\clearfield{annotation}\clearfield{keywords}\clearfield{file}}



\newtheorem{theorem}{Theorem}[section]
\newtheorem{proposition}[theorem]{Proposition}
\newtheorem{definition}[theorem]{Definition}
\theoremstyle{remark}
\newtheorem{remark}[theorem]{Remark}

\newcommand{\cmotion}{\mbox{C-Motion}}
\newcommand{\ElEmEr}{\mbox{ElEmEr}}
\newcommand{\dd}{\mathrm{d}}
\newcommand{\ii}{\mathrm{i}}
\newcommand{\ee}{\mathrm{e}}
\newcommand{\CC}{\mathbb{C}}
\newcommand{\SSph}{\mathbb{S}^{2}}
\newcommand{\SU}{\mathrm{SU}(2)}
\newcommand{\SO}{\mathrm{SO}(3)}
\newcommand{\Tr}{\operatorname{Tr}}
\newcommand{\sgn}{\operatorname{sgn}}
\newcommand{\ahat}{\hat{\bm a}}
\newcommand{\bhat}{\hat{\bm b}}
\newcommand{\nhat}{\hat{\bm n}}
\newcommand{\bsigma}{\bm\sigma}
\newcommand{\epsm}{\varepsilon}
\newcommand{\Jmap}{\mathcal J}
\newcommand{\Crel}{C^{\mathrm{rel}}}
\newcolumntype{P}[1]{>{\raggedright\arraybackslash}p{#1}}

\hypersetup{
  colorlinks=true,
  linkcolor=blue!50!black,
  citecolor=blue!50!black,
  urlcolor=blue!50!black,
  pdftitle={Structural Entanglement in C-Motion I: Common-Source Pairing and the Singlet Preparation Kernel},
  pdfauthor={Toma Kavonius}
}

\title{\textbf{Structural Entanglement in \cmotion{} I:}\\
Common-Source Pairing and the Singlet Preparation Kernel}
\author{Toma Kavonius}
\date{24 November 2026}

\begin{document}
\maketitle

\begin{center}
\small\itshape
Scope and limitations.  This manuscript isolates a bounded common-source preparation
result for singlet-type correlations.  The physical analyzer-transduction law,
native Born weighting, single-trial outcome dynamics, exchange statistics, and
full operator microcausality remain unresolved where stated.
\end{center}
\vspace{0.6em}

\begin{abstract}
Two closures formed by one source may retain a relational structure before
measurement. This paper asks what common-source preparation can establish
independently of analyzer response and outcome selection. In a specified
Conditional C-Motion effective sector, a rotationally scalar parent creates
two positive-frequency rotational doublets through a leading local
complex-bilinear formation vertex. Rotational covariance selects the unique
nontrivial alternating pairing $a^T\epsm b$, with $\epsm=i\sigma_2$, up to
overall strength. Its anti-linear preparation map $\Jmap=\epsm K$,
where $K$ denotes complex conjugation, intertwines
the $\SU$ rotation action while reversing Pauli-vector orientation. In
matched prepared daughter frames this yields the Conditional relational
tensor $\Crel_{ij}=-\delta_{ij}$. It is not a measured covariance, and
physical doublet selection, the microscopic source law, and separated-frame
transport remain unestablished.

Even when isotropy, continuity, signed-axis symmetry, equal-setting
anticorrelation, unbiased marginals and no-signalling are supplied,
two-analyzer correlations need only have the form
$E(\ahat,\bhat)=f(\ahat\cdot\bhat)$; the singlet function $f(x)=-x$
is not selected. An additional physical analyzer-transduction law is needed;
a bilinear law would conditionally recover the singlet angular dependence.
Standard singlet probabilities and Bell/CHSH bounds serve as downstream
controls, not native preparation or measurement dynamics. Native Born
weighting, single-trial outcomes, exchange statistics and full operator
microcausality remain Open.
\end{abstract}


% =====================================================================
\section{The physical question}
\label{sec:physical_question}
% =====================================================================

Consider one preparation event that produces two localized closures and sends them toward two later
measurement stations.  After separation, the closures occupy different regions of physical
three-space, but spatial separation alone does not require their complete physical description to be
the sum of two independent local descriptions.  In the field-first \cmotion{} ontology, both are
organizations of one complete \ElEmEr{} field state, and a common source can prepare a relation
between them before the analyzer settings are chosen.

The proposal is therefore not that one wing sends information to the other during measurement.  The
physical picture is instead
\begin{equation}
  \boxed{\text{one source event prepares one relationally constrained two-closure state}.}
  \label{eq:common_source_picture}
\end{equation}
A relation written at preparation may persist as part of the state carried into the two separated
regions.  Whether that relation has the right form for singlet correlations is a separate question;
common ancestry by itself is not enough.

Three layers must be kept distinct from the start:
\begin{equation}
\boxed{
\begin{array}{ll}
1.&\text{prepared relational structure},\\
2.&\text{local analyzer transduction},\\
3.&\text{single-trial outcome selection}.
\end{array}}
\label{eq:three_layers}
\end{equation}
Preparation asks what relation the common source creates.  Analyzer transduction asks how a chosen
laboratory direction probes that relation.  Outcome selection asks why one particular detector
record occurs in one run.  Structural Entanglement I is primarily about the first layer.  It reaches
the second only far enough to isolate the exact missing law.  The third belongs to the separate measurement-completion problem, which remains scientifically open in \cmotion{} and is outside the scope of this paper.

This ordering also prevents a circular argument.  The paper does not begin by assuming the
ideal singlet law
\begin{equation}
  E(\ahat,\bhat)=-\ahat\cdot\bhat
  \label{eq:target_law}
\end{equation}
and then rename its mathematical structure in \cmotion{} language.  Instead it asks what relation a
specified common-source process can prepare before \eqref{eq:target_law} is supplied.

The central reader question is therefore:
\begin{quote}
\emph{What does \cmotion{} contribute to the singlet problem before the standard quantum measurement
law is supplied?}
\end{quote}
The answer developed below is bounded but nontrivial: within a specified Conditional two-daughter
sector, the source/process symmetry selects an alternating relational pairing and hence an isotropic
negative Pauli-orientation transport relation in matched daughter frames.  The physical
rules for transporting that relation and converting it into exact analyzer statistics remain Open.

% =====================================================================
\section{What is assumed before the preparation calculation}
\label{sec:conditional_sector}
% =====================================================================

The preparation theorem does not start from primitive \ElEmEr{} dynamics alone.  It uses a bounded
effective sector whose physical selection is not yet derived.

\subsection{Two effective rotational components}

Let the two daughter channels be represented by amplitudes
\begin{equation}
  a,b\in\CC^2.
  \label{eq:doublets}
\end{equation}
The symbols \(a\) and \(b\) denote the two-component internal amplitudes assigned to the two daughter
slots.  Under a common lifted spatial rotation they transform together as
\begin{equation}
  a\mapsto Ua,
  \qquad
  b\mapsto Ub,
  \qquad
  U\in\SU.
  \label{eq:doublet_rotation}
\end{equation}
A ``doublet'' here simply means that two effective internal components transform together under this
two-component rotational representation.  The calculation tests what follows \emph{if} the physical
daughter closures occupy such a sector.  A representation-theoretic observation provides
a conditional motivation for testing this sector: if a returned
amplitude must distinguish the two signs of lifted director/frame
$\mathbb Z_2$ holonomy, the minimal readable \(\SU\) carrier has
$j=\frac12$ and is a doublet.  Indeed, the central element $-I\in\SU$
acts in the spin-$j$ irreducible representation as $(-1)^{2j}I$.
Distinguishing its sign requires half-integer $j$; the smallest such
representation has $j=\tfrac12$ and dimension $2j+1=2$.
That representation-theoretic observation
does not establish that actual matter closures require such readability or
select and populate this sector dynamically.  \cmotion{} has not yet
derived from primitive closure dynamics why these daughter modes are realized.  The Pauli matrices below provide a Hermitian basis for this effective rotational
representation; its anti-Hermitian rotation generators are \(T_j=-\ii\sigma_j/2\).
Before the standard-control section, these are not assumed to be laboratory
spin-measurement operators.

\subsection{Why the positive-frequency complex sector matters}

The effective description also selects positive-frequency daughter modes.  Physically, this
specifies one orientation of local phase evolution in the mode decomposition and distinguishes
outgoing daughter content from its conjugate-frequency partner.  Complex notation then packages the
corresponding real quadratures efficiently.  This is more than decorative notation because the
frequency orientation distinguishes the formation process from a conversion process below.

The physical origin of the global complex structure and positive-frequency/positive-energy
selection remains upstream Open/Conditional in the Quantum Foundations companion paper
\cite{CmotionQuantumFoundations2026}.  Nothing in this paper promotes that
selection to a native result.

\subsection{Two useful representation negatives}

Simpler orientation objects do not already provide the needed signed relation.  A headless director
obeys \([d]=\{d,-d\}\) and therefore has the natural nematic representative
\begin{equation}
  Q_{ij}=d_i d_j-\frac13\delta_{ij}.
  \label{eq:nematic}
\end{equation}
Because \(d\) and \(-d\) are physically identified, this representation does not by itself supply a
signed rank-one spin-like relation.

Ordinary spatial inversion is also insufficient.  An axial vector transforms under an orthogonal
map \(O\) as \(\bm s\mapsto(\det O)O\bm s\); for \(O=-I\), the two minus signs cancel.  The negative
relation derived below is therefore not an ordinary parity flip.  These are scoped representation
negatives, not general no-go theorems for closure geometry.

% =====================================================================
\section{The source process selects the relevant pairing class}
\label{sec:process_class}
% =====================================================================

The process-class question ``why \(\mathcal J\) rather than \(I\)?'' sounds arbitrary if the two maps are
presented as competing mathematical choices.  In the bounded effective model they correspond to
different physical process structures.

Let a scalar source mode and the two daughter modes have harmonic dependence
\begin{equation}
\begin{split}
  S(s)&=S_0\ee^{-\ii\Omega_S s},\\
  f_A(s)&=a\,\ee^{-\ii\Omega_A s},\\
  f_B(s)&=b\,\ee^{-\ii\Omega_B s}.
\end{split}
\label{eq:harmonic_modes}
\end{equation}
Here \(s\) is the succession/evolution parameter used by the effective description.  It is not
introduced as a fourth spatial-like dimension or as a directly readable universal clock.

A conversion or mixing process acts on already existing mode content.  Its natural invariant is
sesquilinear,
\begin{equation}
  f_A(s)^\dagger f_B(s),
  \qquad
  \Omega_S\simeq\Omega_A-\Omega_B,
  \label{eq:conversion_channel}
\end{equation}
so a corresponding source coupling \(S(s)f_A(s)^\dagger f_B(s)\) belongs to a
difference-frequency channel.  The matrix product \(a^\dagger b\) is its constant-amplitude
factor, not the complete oscillating mode expression.

The preparation studied here is instead a genuine two-daughter formation process,
\begin{equation}
  S\longrightarrow A+B,
  \qquad
  \Omega_S\simeq\Omega_A+\Omega_B.
  \label{eq:formation_channel}
\end{equation}
With two outgoing positive-frequency daughter slots, the leading local derivative-free scalar
depends complex-bilinearly on the daughters.  The mathematical question is therefore not ``which
pairing does nature arbitrarily prefer?'' but rather:
\begin{quote}
Which complex-bilinear pairing of two daughter doublets remains invariant when both are rotated
together?
\end{quote}
This process distinction is Conditional on the effective mode sector and does not imply that every
physical pair-creation process in nature must have the same form.

% =====================================================================
\section{The alternating invariant pairing}
\label{sec:pairing}
% =====================================================================

\subsection{The compact invariance calculation}

Write the most general local complex-bilinear scalar as
\begin{equation}
  B(a,b)=a^T M b,
  \label{eq:general_bilinear}
\end{equation}
where \(M\) is a constant \(2\times2\) pairing matrix.  Rotational scalarity means that rotating both
daughters together does not change this source coupling:
\begin{equation}
  B(Ua,Ub)=B(a,b)
  \quad\Longleftrightarrow\quad
  \boxed{U^T M U=M\quad\forall U\in\SU.}
  \label{eq:invariance}
\end{equation}

\begin{proposition}[Conditional two-daughter pairing result]
Within the rotationally scalar-parent, positive-frequency fundamental-doublet, local leading
derivative-free complex-bilinear formation class, the invariant space satisfying
\eqref{eq:invariance} is one-dimensional.  Hence
\begin{equation}
  \boxed{
  M=\lambda\epsm,
  \qquad
  \epsm=i\sigma_2=
  \begin{pmatrix}
  0&1\\[-1mm]
  -1&0
  \end{pmatrix},}
  \label{eq:epsilon}
\end{equation}
with overall strength \(\lambda\).
\end{proposition}

\begin{proof}
Write \(M=(m_{ij})\).  For the diagonal subgroup
\[
  U_z(\phi)=\operatorname{diag}
  \!\left(\ee^{-\ii\phi/2},\ee^{\ii\phi/2}\right),
\]
condition \eqref{eq:invariance} for every \(\phi\) forces \(m_{11}=m_{22}=0\).  Applying a
noncommuting rotation, for example one generated by \(\sigma_1\), then gives
\(m_{21}=-m_{12}\).  Therefore \(M\) is proportional to the alternating matrix \(\epsm\).  Direct
substitution verifies \(U^T\epsm U=\epsm\) for every \(U\in\SU\).
\end{proof}

The matrix \(\epsm\) is simply the invariant antisymmetric pairing form of a two-component rotational
doublet.  The notation \(i\sigma_2\) is a convenient standard representation; its appearance does
not independently derive physical electron spin, the physical doublet sector, or fermionic exchange
statistics.  Pair-slot alternation, \(B(b,a)=-B(a,b)\), does not
by itself establish exchange statistics for two identical completed closures.
Nor does it require antisymmetry of a later measured correlation under
interchange of analyzer sites.  In particular, selecting an
\emph{invariant source bilinear} does not prove that
native \ElEmEr{} dynamics prepares a normalized, pure two-particle singlet state.  Such an
identification needs a physical daughter-state composition rule, a source evolution that
populates the relevant two-slot sector, and an appropriate normalization/statistical
interpretation.  These are not supplied by the uniqueness proof.

This resolves the process-class ``\(\mathcal J\) or \(I\)?'' ambiguity only within the stated process
model.  The identity-like sesquilinear structure belongs naturally to a difference-frequency
conversion channel, whereas the genuine scalar-parent two-daughter formation channel selects the
alternating bilinear.  The result is therefore best described as a \emph{\cmotion{}-constrained
Conditional preparation bridge}, not as an N1 derivation from primitive \ElEmEr{} dynamics.

\subsection{The anti-linear preparation relation}

A real interaction carrying the selected pair channel can be represented schematically
at the local formation vertex by the \emph{full mode fields} of
\eqref{eq:harmonic_modes},
\begin{equation}
  V_{\rm pair}(s)
  =-g_S\!\left[
  S(s)\bigl(f_A(s)^T\epsm f_B(s)\bigr)^*
  +S(s)^*\bigl(f_A(s)^T\epsm f_B(s)\bigr)
  \right],
  \label{eq:pair_interaction}
\end{equation}
where \(g_S\) denotes an effective source coupling.  The first term contains
\begin{equation}
  S(s)\bigl(f_A(s)^T\epsm f_B(s)\bigr)^*
  =S_0\,(a^T\epsm b)^*\,
    \ee^{-\ii(\Omega_S-\Omega_A-\Omega_B)s},
  \label{eq:pair_phase}
\end{equation}
so a slowly varying (resonant) channel is possible for
\(\Omega_S\simeq\Omega_A+\Omega_B\) as in \eqref{eq:formation_channel}.
Variation of the conjugated pair factor with respect to the daughter conjugate
fields produces terms proportional to \(\epsm f_B^*\) and
\(-\epsm f_A^*\), respectively; the relative minus sign is the
antisymmetry \(\epsm^T=-\epsm\).  The constant amplitudes \(a,b\) alone
would not show the resonance phase.  This motivates the preparation relation
\begin{equation}
  \boxed{\Jmap=\epsm K.}
  \label{eq:Jmap}
\end{equation}
Up to an overall phase, this is the algebraic form of Wigner's
spin-$\tfrac12$ time-reversal operator \cite{Wigner1932TimeReversal}:
$\Jmap^2=-I$ and the Pauli vector changes sign.  This algebraic
identification does not make the source preparation a physical reversal
of time, momentum, or an entire field history.
Here \(K\) means complex conjugation.  Because complex conjugation reverses multiplication by
\(i\), \(\Jmap\) is anti-linear rather than linear.

Physically, \(\Jmap\) maps the effective internal description in one daughter slot to the
complementary relational orientation assigned to the partner by the common preparation.  It is a
relation between the prepared daughter internal spaces, not a signal propagating continuously from
one separated wing to the other.  It is also not identified here with charge conjugation or with the
exchange of two identical completed closures.

% =====================================================================
\section{From the pairing map to the relational tensor}
\label{sec:relational_tensor}
% =====================================================================

The alternating pairing fixes a \emph{dual Pauli-orientation relation} between
the daughter slots in the chosen common rotational frame.  Its matrix identity is
\begin{equation}
  \epsm\sigma_j^T\epsm^{-1}=-\sigma_j,
  \qquad j=1,2,3.
  \label{eq:generator_transport}
\end{equation}
Define \(\Crel_{ij}\) in these matched bases by
\begin{equation}
  \epsm\sigma_j^T\epsm^{-1}
  =\Crel_{ij}\sigma_i.
  \label{eq:Crel_definition}
\end{equation}
Comparison with \eqref{eq:generator_transport} gives
\begin{equation}
  \boxed{\Crel_{ij}=-\delta_{ij}.}
  \label{eq:Crel}
\end{equation}

\begin{definition}[Relational Pauli-orientation transport tensor]
\(\Crel_{ij}\) records the dual Pauli-basis orientation relation induced
by the alternating common-source pairing in \emph{matched daughter frames}.
It is Conditional structural information about the preparation, not a
measured two-analyzer correlation tensor.
\end{definition}

The sign must not be misread as a reversal of actual \(\SU\) rotations.
Since \(\Jmap=\epsm K\) is anti-linear, with \(K\) complex conjugation,
\begin{equation}
  \Jmap\sigma_j\Jmap^{-1}=-\sigma_j,
  \qquad
  T_j:=-\frac{\ii}{2}\sigma_j,
  \qquad
  \boxed{\Jmap T_j\Jmap^{-1}=T_j},
  \label{eq:antilinear_rotation_guard}
\end{equation}
and hence \(\Jmap U\Jmap^{-1}=U\) for \(U\in\SU\).
Anti-linearity reverses \(\ii\) as well as the Pauli-vector sign: the map
\emph{intertwines} the \(\SU\) representation; it does not reverse
its anti-Hermitian rotation generators.  Thus \(\Crel=-I\) describes
opposite dual Pauli orientation rather than reversal of physical rotational
evolution.

\paragraph{Separated daughter frames.}
The form \(\Crel=-I\) is stated in the matched reference frames fixed by
preparation.  To compare differently oriented daughter frames, suppose their
Pauli-axis component bases are carried by orthogonal rotations
\(R_A,R_B\in\SO\) into a common comparison convention.  For this
explicit basis convention, the same transported relation has components
\begin{equation}
  C^{\mathrm{rel,comp}}=-R_A R_B^T.
  \label{eq:frame_transport}
\end{equation}
This is conditional frame bookkeeping, not an independently derived
physical parallel-transport or interaction law.  Purely changing a basis
also changes the components of each analyzer axis and cannot alter a
physical scalar correlation.  Actual separated closures may undergo
different physical rotations, couplings or loss of preparation coherence;
how their frames and relation evolve is an Open input from native dynamics.
No preferred global frame or instantaneous transmission between the
separated wings is assumed.

Equation \eqref{eq:Crel} therefore states an opposite dual Pauli
orientation in the prepared matching convention.  It does \emph{not} yet
say what detector outcomes occur when an analyzer is applied.  In particular,
\begin{equation}
  \boxed{
  \Crel_{ij}=-\delta_{ij}
  \quad\text{is not yet}\quad
  \langle A_iB_j\rangle=-\delta_{ij}.}
  \label{eq:rel_not_cov}
\end{equation}
The left side belongs to preparation.  The right side would require a definition of local measured
observables, an analyzer-response law, and a statistical rule for outcomes.

A classical comparison makes the distinction concrete.  If an ordinary hidden unit vector
\(\nhat\) were isotropically distributed and the partner carried \(-\nhat\), then
\begin{equation}
  \left\langle n_i(-n_j)\right\rangle
  =-\int_{\SSph}n_i n_j\frac{\dd\Omega}{4\pi}
  =-\frac13\delta_{ij}.
  \label{eq:classical_covariance}
\end{equation}
Thus the \(-\delta_{ij}\) in \eqref{eq:Crel} cannot simply be reinterpreted as the covariance of an
ordinary isotropic hidden spatial vector.  Its meaning is dual Pauli-orientation transport in matched frames, not a
classical ensemble average.

% =====================================================================
\section{The analyzer-transduction boundary}
\label{sec:analyzer}
% =====================================================================

The next question is how a local laboratory analyzer turns the prepared relation into a measured
two-site statistic.  Let \(\ahat\) and \(\bhat\) be unit directed analyzer axes, let
\(A,B\in\{+1,-1\}\) be the recorded signs, and write \(E(\ahat,\bhat)=\langle AB\rangle\).

\subsection{Isotropy leaves an arbitrary angular function}

If the complete preparation has no preferred direction, a simultaneous rotation of both analyzer
axes cannot change the scalar correlation:
\begin{equation}
  E(R\ahat,R\bhat)=E(\ahat,\bhat),
  \qquad R\in\SO.
  \label{eq:rotational_E}
\end{equation}
For two directed unit vectors in three dimensions, simultaneous rotational invariance leaves only their inner product, so
\begin{equation}
  \boxed{E(\ahat,\bhat)=f(\ahat\cdot\bhat)}
  \label{eq:general_f}
\end{equation}
for some bounded function \(f:[-1,1]\to[-1,1]\).  Isotropy therefore fixes the variable on which the
correlation may depend, but not the function of that variable.

Suppose in addition that reversing a directed analyzer reverses the signed local outcome.  Then
\(f(-x)=-f(x)\).  Supplying perfect equal-setting anticorrelation gives \(f(1)=-1\) and
\(f(-1)=+1\).  Even with continuity, many functions remain.  Two examples are
\begin{equation}
  f(x)=-x^3,
  \qquad
  f(x)=-\sin\!\left(\frac{\pi x}{2}\right).
  \label{eq:counterexamples}
\end{equation}
Both satisfy the same symmetry and endpoint conditions as the singlet law.

\paragraph{Conditional symmetry route to unbiased local marginals.}
There is a narrower way to obtain \emph{unbiased} one-site marginals
without postulating them independently.  Suppose that (i) the mean
\(m_A(\ahat)=\langle A\rangle\) depends only on the local analyzer setting
and the relevant reduced local preparation/response data, not on the remote
setting; (ii) that \emph{complete} reduced local comparison is rotationally
neutral, so \(m_A(R\ahat)=m_A(\ahat)\) for all \(R\in\SO\); and (iii)
reversing the directed analyzer reverses its signed mean.  A proper
rotation maps \(\ahat\) to \(-\ahat\), and therefore
\begin{equation}
  m_A(-\ahat)=m_A(\ahat)=-m_A(\ahat)
  \quad\Longrightarrow\quad
  \boxed{m_A(\ahat)=0,}
  \label{eq:conditional_local_marginal}
\end{equation}
and likewise \(m_B(\bhat)=0\) under the analogous premises.  For binary
outcomes this gives \(P(A=\pm1)=P(B=\pm1)=\tfrac12\).
This is a \emph{Conditional} response-symmetry lemma, not a native
no-signalling or Born theorem: remote-setting independence is an explicit
premise.  A scalar-parent label by itself does not establish (ii).  For
example, conditioning on daughter flight directions, anisotropic
apparatus states or other selected geometry can leave physical local axes
and invalidate the required full reduced-state isotropy.

If unbiased binary marginals are supplied, or obtained through the
separately stipulated premises above, every bounded \(f\) generates the
pairwise joint family
\begin{equation}
  P(A,B\mid\ahat,\bhat)
  =\frac14\left[1+AB\,f(\ahat\cdot\bhat)\right],
  \label{eq:general_joint}
\end{equation}
whose one-site marginals are \(1/2\), independent of the distant setting.  Marginal no-signalling
therefore does not choose \(f(x)=-x\) either.  The logical insufficiency
of no-signalling to select quantum correlations was established by
Popescu and Rohrlich \cite{PopescuRohrlich1994}.  The continuous,
rotationally invariant families used here are the present construction;
they are not identified with the finite-setting PR-box law.

\begin{proposition}[Analyzer-law underdetermination]
Rotational isotropy, continuity, signed-axis oddness, perfect equal-setting anticorrelation, unbiased
binary marginals, and marginal no-signalling do not uniquely determine the singlet angular law.
\end{proposition}

The counterexamples are pairwise conditional probability laws, not automatically
joint distributions over outcomes for every possible analyzer setting, and they
need not be realizable by a quantum state.  To make the distinction quantitative,
take the second function in \eqref{eq:counterexamples} and the usual coplanar
CHSH settings with
\(\ahat\cdot\bhat=\ahat\cdot\bhat'=\ahat'\cdot\bhat=1/\sqrt2\)
and \(\ahat'\cdot\bhat'=-1/\sqrt2\).  Defining
\(S=E(\ahat,\bhat)+E(\ahat,\bhat')+E(\ahat',\bhat)-E(\ahat',\bhat')\),
one obtains
\begin{equation}
  |S|=4\sin\!\left(\frac{\pi}{2\sqrt2}\right)
  \simeq3.58>2\sqrt2.
  \label{eq:nonquantum_nosignal_chsh}
\end{equation}
Thus even a rotationally invariant, continuous, unbiased, no-signalling
pairwise law with perfect equal-setting anticorrelation can exceed the
standard quantum Tsirelson bound \cite{Tsirelson1980}.  The example is a \emph{mathematical
nonquantum control}, not a physical prediction of \cmotion{}.  Matching the
observed singlet requires selecting the \emph{quantum-admissible} correlation,
not simply obtaining no-signalling or a Bell violation.

This underdetermination is a scientific result, not an editorial caveat.  It shows that the exact
singlet angular dependence contains constitutive information about how an analyzer couples to the
prepared relation.

\subsection{The missing bilinear law}

The narrow missing bridge can be written as
\begin{equation}
  \boxed{
  E(\ahat,\bhat)
  =g\,\ahat_i\Crel_{ij}\bhat_j.}
  \label{eq:analyzer_law}
\end{equation}
Here $\ahat_i$ and $\bhat_j$ are Cartesian components of the directed analyzer
axes $\ahat$ and $\bhat$, and $g$ is a possible response normalization.
Physically, \eqref{eq:analyzer_law} says that each local analyzer interrogates the
prepared rotational relation linearly in its chosen direction, so the two-site mean
is the double projection of the relational tensor.

If a future physical derivation supplies this law, and the two analyzer axes
are expressed in the matched prepared frames (or consistently referred to a
common transported comparison frame), \eqref{eq:Crel} gives
\begin{equation}
  E(\ahat,\bhat)=-g\,\ahat\cdot\bhat.
  \label{eq:dot_up_to_g}
\end{equation}
For independently reoriented frames, \eqref{eq:frame_transport} must
instead be used with consistently transformed analyzer directions.
The physical frame transport and analyzer comparison remain Open.
Perfect equal-setting anticorrelation would then fix only the residual normalization, \(g=1\).  The
intended architecture would therefore be
\begin{equation}
  \boxed{\text{derive the form}\;\longrightarrow\;\text{calibrate the residual normalization}.}
  \label{eq:form_then_norm}
\end{equation}
At present even this architecture is only Conditional because exact analyzer bilinearity has not
been derived from a physical closure--apparatus interaction.

A weak-response expansion gives a useful clue but not a completion.  If the
respective local analyzer couplings begin as
\begin{equation}
  V_A\sim-h_A\,\ahat_iG_i^{(A)},
  \qquad
  V_B\sim-h_B\,\bhat_jG_j^{(B)},
  \label{eq:local_weak_analyzer_couplings}
\end{equation}
then the leading connected cross-response naturally has the form
\begin{equation}
  \delta^2\mathcal R_{AB}
  \propto
  \ahat_i\chi_{ij}\bhat_j.
  \label{eq:weak_response}
\end{equation}
Thus ordinary perturbative response produces a bilinear first harmonic.  Exact detector statistics
may nevertheless become nonlinear through thresholding, saturation, record formation, or other
nonperturbative response.  Equation \eqref{eq:weak_response} is therefore a future-work clue, not a
derivation of \eqref{eq:analyzer_law}.

% =====================================================================
\section{Why ordinary shared-vector models are not enough}
\label{sec:shared_vector}
% =====================================================================

Two concise shared-vector controls separate the relational preparation structure
from an ordinary shared spatial direction.  Their calculated correlations are
valid within the supplied classical models; the Negative conclusion concerns
only their identification with the full singlet angular law.

For a uniformly distributed hidden direction \(\nhat\), consider deterministic local threshold
readout
\begin{equation}
  A=\sgn(\ahat\cdot\nhat),
  \qquad
  B=-\sgn(\bhat\cdot\nhat).
  \label{eq:threshold_model}
\end{equation}
This Bell-type model \cite{Bell1964,CHSH1969} gives, for
\(\theta=\arccos(\ahat\cdot\bhat)\),
\begin{equation}
  E_{\rm th}(\theta)=-1+\frac{2\theta}{\pi},
  \label{eq:threshold_result}
\end{equation}
not \(-\cos\theta\).

A smooth factorized local response also misses the target.  If
\(\langle A\rangle=\ahat\cdot\nhat\),
\(\langle B\rangle=-\bhat\cdot\nhat\), and the wings are conditionally independent given
\(\nhat\), then
\begin{equation}
  E_{\rm fact}(\ahat,\bhat)
  =-\int_{\SSph}(\ahat\cdot\nhat)(\bhat\cdot\nhat)
  \frac{\dd\Omega}{4\pi}
  =-\frac13\ahat\cdot\bhat.
  \label{eq:factorized_result}
\end{equation}
The same \(1/3\) factor appears in the classical covariance \eqref{eq:classical_covariance}; see also
standard analyses of isotropic vector hidden-variable models \cite{GhirardiRomano2012}.  These
failures do not prove the \cmotion{} preparation mechanism.  They show only that common ancestry plus
an ordinary isotropic hidden direction is structurally insufficient.

% =====================================================================
\section{Downstream standard singlet controls}
\label{sec:standard_control}
% =====================================================================

We now switch from the \cmotion{} preparation question to the standard effective quantum
representation.  This section states what follows if the missing analyzer bridge reproduces the
usual singlet measurement rule.  None of the following mathematics is used to select \(\epsm\) or to
derive \(\Crel\).

In the \emph{supplied standard quantum two-doublet model}, the same invariant
antisymmetric tensor is represented by the normalized singlet ket
\begin{equation}
  |\Psi^-\rangle
  =\frac{1}{\sqrt2}
  \left(|+\rangle_A|-\rangle_B-|-\rangle_A|+\rangle_B\right).
  \label{eq:singlet_state}
\end{equation}
This is a correspondence between invariant tensors, \emph{not} a derivation
that the microscopic two-closure source prepares this pure Hilbert-space state.
The tensor-product composition rule, amplitude normalization, and physical
state-production dynamics are supplied or remain Open; if the daughters are
identical, the required combined internal/spatial exchange structure is
separately unresolved.  With the standard Pauli analyzer rule and Born trace,
this \emph{assumed effective state} gives
\begin{equation}
  E(\ahat,\bhat)=-\ahat\cdot\bhat.
  \label{eq:standard_E}
\end{equation}
In matched site/analyzer conventions this standard correlation is symmetric
under interchange of its two axes, \(E(\ahat,\bhat)=E(\bhat,\ahat)\),
despite the antisymmetry of the preparation pairing slots.  This
distinction does not decide physical exchange statistics.
The corresponding binary joint distribution is
\begin{equation}
  \boxed{
  P(A,B\mid\ahat,\bhat)
  =\frac14\left(1-AB\,\ahat\cdot\bhat\right).}
  \label{eq:born_joint}
\end{equation}
Equation \eqref{eq:born_joint} follows from the supplied standard singlet/Born description.  It is not
derived directly from \(\Crel\).

Summing over the distant outcome gives
\begin{equation}
  P(A\mid\ahat,\bhat)=\frac12,
  \qquad
  P(B\mid\ahat,\bhat)=\frac12,
  \label{eq:marginals}
\end{equation}
so changing the distant analyzer setting cannot change the local marginal distribution.  This is
marginal no-signalling.  It does not constitute a full microscopic locality theorem.

For the standard CHSH analyzer arrangement \cite{CHSH1969,Tsirelson1980},
\eqref{eq:standard_E} gives
\begin{equation}
  |S|=2\sqrt2.
  \label{eq:CHSH}
\end{equation}
Equation \eqref{eq:CHSH} is the ideal prediction for the supplied state and
optimal analyzer settings, not an experimentally established exact saturation.
Loophole-free experiments report finite-precision Bell violations with
finite visibility and implementation-dependent statistics.  Hensen et al.
test electron spins \cite{Hensen2015LoopholeFree}; Giustina et al. and
Shalm et al. test photonic Bell inequalities
\cite{Giustina2015LoopholeFree,Shalm2015LoopholeFree}.
These experiments test Bell constraints; they need not implement this
exact ideal singlet/CHSH arrangement.  Bell/CHSH therefore checks the
nonclassical correlation structure that a completed bridge must reproduce.  Bell violation is not evidence specific to the preparation
mechanism proposed here, and this paper does not derive the Tsirelson bound from a native \cmotion{}
operator algebra.

% =====================================================================
\section{Local causality is not separable completeness}
\label{sec:locality}
% =====================================================================

The preparation picture distinguishes two propositions that are often blurred together.

\emph{Local causal propagation} concerns how newly generated physical changes propagate.  The Main
Framework \cite{CmotionMainFramework2026} constrains that propagation by the
framework's local-\(c\) causal architecture.  By
contrast, \emph{separable completeness} is the stronger claim that the complete physical state of a
prepared pair is nothing more than two independently complete local states, one at each wing.

A common-source relational state can reject the second proposition without asserting a
superluminal signal during measurement.  The distinction is
\begin{equation}
  \boxed{
  \text{local causal propagation}
  \;\neq\;
  \text{separable value-completeness}.}
  \label{eq:locality_distinction}
\end{equation}
For a proposed \emph{complete} shared preparation variable \(\lambda\), Bell-local
factorization reads
\begin{equation}
  P(A,B\mid\ahat,\bhat,\lambda)
  =P(A\mid\ahat,\lambda)P(B\mid\bhat,\lambda).
  \label{eq:bell_factorization}
\end{equation}
Together with setting-independent preparation weights
\(\rho(\lambda\mid\ahat,\bhat)=\rho(\lambda)\), ordinary normalized
conditional probabilities and binary outcomes, this yields the CHSH bound \cite{Bell1964,CHSH1969}
\begin{equation}
  \boxed{|S|\leq2,}
  \qquad
  S:=E(\ahat,\bhat)+E(\ahat,\bhat')
  +E(\ahat',\bhat)-E(\ahat',\bhat').
  \label{eq:bell_complete_bound}
\end{equation}
Here \(\lambda\) may describe the \emph{entire jointly prepared pair},
including an arbitrarily elaborate nonseparable/common-source relation;
it need not be a pair of independent hidden directions.  The inequality
still follows if \(\lambda\) screens off the two local response laws and
is independent of the later setting choices.  The ideal singlet prediction \(|S|=2\sqrt2\), and measured Bell
violations above the relevant local bound, cannot be reproduced under all
those premises.  Exact experimental saturation of the ideal value is not
required for this exclusion.

Local finite-speed propagation and operational no-signalling alone do not
logically establish Bell factorization.  Conversely, claiming a shared
nonseparable preparation is \emph{not by itself} an escape from the Bell
constraint: a completed physical analyzer/outcome model must say precisely
whether it satisfies the screening-off, setting-independence and related
premises, and demonstrate the actual joint statistics.  This manuscript
neither selects a Bell-violating mechanism nor provides such a completed
model.  The problem of physical analyzer coupling, probabilities and outcomes
remains Open, with measurement completion owned downstream.

The conceptual distinction does not by itself derive Bell correlations.
Moreover,
\begin{equation}
  \boxed{\text{marginal no-signalling}\;\neq\;\text{full operator microcausality}.}
  \label{eq:no_signal_not_micro}
\end{equation}
The standard joint law has no-signalling marginals, while spacelike commutation of a future complete
extended-mode operator theory remains an Open proof obligation.

% =====================================================================
\section{Scientific boundary and compact audit register}
\label{sec:boundary}
% =====================================================================

The preparation result narrows the remaining problem rather than completing measurement.  The
immediate missing bridge is
\begin{equation}
  \boxed{
  \text{prepared relational structure}
  \;\longrightarrow\;
  \text{analyzer-dependent statistical correlation}.}
  \label{eq:sharp_boundary}
\end{equation}
The microscopic source/action that realizes the selected process class also remains unresolved, but
it is upstream of the bounded theorem rather than silently assumed to be solved.

The principal statuses are collected here so that the main argument can remain organized by physics.

The register below records the scientific status of the claims used in
this manuscript.  In the Derivation Ledger's vocabulary, Conditional means
that a conclusion follows from an explicit premise or model class not itself
derived; Structural means an exact or banked organizational/accounting statement
that does not supply a microscopic law; Model-derived means rigorously derived
inside a specified effective or analogue model, never native Derived.
Open means unresolved; Open (sharp) identifies a sharply isolated missing
bridge, not greater confidence.  Negative excludes an identification, mechanism
or branch only within its stated scope.  Publication deferral does not resolve
an Open scientific question.

Scientific status and provenance are separate.  N1 denotes a primitive/native
route and N2 a C-Motion-constrained effective bridge; standard/control denotes
supplied conventional mathematics.  These route labels do not constitute a
second scientific-status vocabulary.  No result here is promoted to N1.

\begingroup
\small
\setlength{\tabcolsep}{3pt}
\begin{longtable}{P{5.0cm}P{3.6cm}P{6.2cm}}
\toprule
\textbf{Claim} & \textbf{Status} & \textbf{Scope / boundary}\\
\midrule
\endfirsthead
\toprule
\textbf{Claim} & \textbf{Status} & \textbf{Scope / boundary}\\
\midrule
\endhead
Common-source relational preparation interpretation & Conditional & Conditional \cmotion{} interpretation inherited from the field-first ontology; this does not select the singlet law.\\
Fundamental doublet daughter sector & Conditional & Two-component lifted rotational representation is the tested effective sector, not a primitive matter derivation; physical sector selection remains Open.\\
Physical complex / positive-frequency sector & Conditional & Used conditionally in the process-class distinction; physical global complex/positive-frequency selection remains Open (sharp).\\
Alternating scalar two-daughter pairing \(a^T\epsm b\) & Conditional & Conditional exact result within the stated N2 scalar-parent, positive-frequency-doublet, local leading complex-bilinear formation class.\\
\(\Jmap=\epsm K\) & Conditional & Conditional N2 anti-linear preparation relation: reverses Pauli-vector signs while intertwining full \(\SU\) rotations; not charge conjugation, communication, or exchange.\\
\(\Crel_{ij}=-\delta_{ij}\) & Conditional & Conditional N2 dual Pauli-orientation relation in matched prepared frames; separated-frame transport law is Open; not measured covariance.\\
Identification \(\Crel_{ij}=\langle A_iB_j\rangle\) & Open & Open bridge requiring local observables, analyzer transduction, and outcome statistics.\\
Rotational reduction \(E=f(\ahat\cdot\bhat)\) & Structural & Exact Structural symmetry consequence under the stated isotropy; it does not select the function \(f\).\\
Isotropy + continuity + signed-axis symmetry + equal-setting anticorrelation + unbiased marginals + no-signalling \(\Rightarrow f=-x\) & Negative & Explicit pairwise counterexamples; one exceeds the quantum Tsirelson bound while preserving no-signalling.\\
Exact bilinear analyzer law
\(E=g\,\ahat_i\Crel_{ij}\bhat_j\) &
Open &
Minimum missing constitutive analyzer law.\\
Native measured law \(E=-\ahat\cdot\bhat\) & Open & Follows only after the analyzer law plus normalization are supplied.\\
Standard singlet/Born representation & Conditional & Supplied tensor-product, pure-singlet and Born control; invariant pairing alone does not derive physical singlet-state production or exchange statistics.\\
Marginal no-signalling & Model-derived & Model-derived/control consequence of the supplied joint law; exact for that law, but not full microscopic locality.\\
Bell/CHSH result & Model-derived & Standard/control consequence of supplied quantum singlet; factorizing complete shared-state models with independent settings obey \(|S|\leq2\), irrespective of relational complexity.\\
Conditional unbiased local marginals & Conditional & Zero one-site means follow only given a rotationally neutral complete reduced local comparison, local-setting-only response, and signed-axis reversal; no native no-signalling theorem.\\
Definite outcomes / Born dynamics / detector records & Open & Scientifically Open; measurement completion is outside the scope of this paper and remains a deferred research problem. \\
Fermionic exchange statistics & Open (sharp) & The sharply isolated exchange-sign bridge remains unresolved: alternating pair-slot pairing and one-body \(4\pi\) return do not derive many-body statistics.\\
Headless director and parity identifications & Negative & Scoped representation controls: a nematic director lacks signed rank-one orientation, and ordinary inversion does not reverse an axial vector; no general geometric no-go.\\
Classical shared-vector correlations & Model-derived & Exact threshold and factorized model results in \cref{sec:shared_vector}; their identification with the full singlet law is Negative, not the calculations themselves.\\
Separated-frame comparison & Conditional & Equation \eqref{eq:frame_transport} is basis bookkeeping with supplied rotations; physical transport and coherence dynamics remain Open.\\
Weak-response bilinearity & Conditional & Leading connected perturbative response with supplied local couplings; exact analyzer statistics and nonlinear detector response remain Open.\\
Full operator microcausality & Open & No-signalling of one-site marginals does not prove spacelike commutation of the full extended-mode operator algebra.\\
\bottomrule
\end{longtable}
\endgroup

\paragraph{Scope of the no-signalling controls.}
The \emph{Model-derived} no-signalling entry in this paper refers specifically
to the exact marginal calculation from the \emph{supplied} standard singlet/Born
joint law \eqref{eq:born_joint}.  In a separate stipulated allowed joint-kernel
class with additional stipulated admissibility and marginal assumptions,
a no-signalling result can be established as an exact conditional
mathematical statement.  This separate control is not used in the preparation
theorem.  Neither statement derives physical analyzer probabilities or full
operator microcausality.  Actual no-signalling for a future native analyzer theory
must be checked from its completed joint law.

The \emph{Conditional} standard singlet/Born control consists of imported
mathematical consequences, not native \cmotion{} Born or outcome dynamics.  These remain scientifically Open in the separate
measurement-completion work (publication-deferred).  The deferred measurement-completion programme's target is not merely
to preserve a postulated probability distribution: it must justify a
physical apparatus/environment sampling measure \(\nu\) and single-trial
selection dynamics \(\mathcal S\) that reproduce the standard frequencies
\(\operatorname{Tr}(\rho E_i)\) for specified preparation and measurement
contexts.  Equivariance could preserve \emph{given} weights but would not,
without a further argument, select those weights physically.

The empirical inputs are also kept separate.  Source isotropy may be an empirical preparation
property if the future microscopic source law does not derive it.  The
conditional local-marginal lemma \eqref{eq:conditional_local_marginal}
requires stronger reduced-response assumptions than a scalar parent alone.
Unbiased marginals are not needed to derive \(\Crel\).  Perfect equal-setting anticorrelation may normalize \(g\) only after the Open
bilinear analyzer law is established.  The full angular law \(E=-\ahat\cdot\bhat\) is never used as
an input to the preparation theorem.

Finally, the representation boundary is strict.  The effective doublet is represented in
\(\CC^2\); \(\epsm\) is its invariant alternating bilinear form; \(\Jmap\) is a preparation relation;
and \(\Crel\) is a dual Pauli-orientation transport tensor in matched frames.  Only in \cref{sec:standard_control} is the familiar
tensor-product singlet, Pauli analyzer algebra, and Born rule supplied as standard effective quantum
machinery.

\paragraph{Future-model admission controls.}
A proposed native analyzer/measurement law must preserve normalization,
positivity, operational no-signalling and membership of its claimed
quantum-admissible joint-kernel class under the \emph{actual} evolution;
agreement with one CHSH number is insufficient.  It must reproduce the full
angular dependence and test stability under source/analyzer variations.
Consistency of multi-pair preparation, composition and swapping is a
separate future test, not a preparation theorem of this paper.  These
controls state what a later physical model must demonstrate, rather than
treating results from prescribed-kernel models as evidence of a native microscopic derivation.

% =====================================================================
\section{Conclusion}
\label{sec:conclusion}
% =====================================================================

The paper asked what a common source can prepare before the standard quantum measurement law is
supplied.  Within the stated Conditional effective sector, the answer is sharper than an unspecified
``shared geometry.''  A rotationally scalar source forming two positive-frequency fundamental
doublet daughters through a leading local complex-bilinear formation process admits the unique
alternating pairing
\begin{equation}
  a^T\epsm b,
  \qquad
  \Jmap=\epsm K,
  \qquad
  \Crel_{ij}=-\delta_{ij}.
  \label{eq:conclusion_result}
\end{equation}
The first expression is the invariant pair-formation relation (whose full
mode-field vertex has the sum-frequency phase shown in \eqref{eq:pair_phase}),
the second is its anti-linear
preparation map, and the third records opposite dual Pauli-axis orientation in matched
prepared daughter frames.  The anti-linear map intertwines physical \(\SU\) rotations;
separated-frame dynamics remain Open.  These are Conditional N2 results, not primitive
N1 derivations.

The result stops at a physically meaningful boundary.  \(\Crel\) is not the measured covariance.
Rotational isotropy gives only
\begin{equation}
  E(\ahat,\bhat)=f(\ahat\cdot\bhat),
  \label{eq:conclusion_general_f}
\end{equation}
and strong symmetry, endpoint, marginal, and no-signalling conditions still do not force
\(f(x)=-x\).  The exact local law that converts the prepared rotational relation into analyzer
statistics remains Open.  If that law is eventually derived as a bilinear double projection, then
\(\Crel=-I\) fixes the angular dependence and equal-setting anticorrelation can fix only the residual
normalization.

The remaining task is therefore narrower than ``derive entanglement from scratch'': determine how a
physical analyzer interrogates the prepared relational structure, after its daughter-frame
transport has been specified.  Both transport and analyzer coupling remain Open.
Standard singlet probabilities, Bell/CHSH violation, and no-signalling marginals show what the
completed bridge must reproduce, but they are downstream controls.  The
invariant pairing does not alone produce a pure physical singlet or bypass
Bell-local factorization of complete shared states; additional analyzer
physics must select the quantum-admissible law.  Born weighting, definite
single-trial outcomes, detector records, physical doublet and complex-sector selection, the
microscopic source/action, exchange statistics, and full operator microcausality remain explicit
open problems rather than hidden assumptions of Structural Entanglement I.



\newpage
\printbibliography

\end{document}